The two parametrisations shipped with these notes have the same branching ratio
$\branchingRatio$ and the same total rate $\totalRate$,
and differ in one structural respect.
The pressure vector $\pressure$ partitions the six types into those that push the mid-price
up and those that push it down.
The first kernel excites \emph{across} that partition:
what depletes a side calls forth what refills it.
The second excites \emph{within} it:
what pushes the price one way calls forth more of the same.
Contemporaneously the two are indistinguishable,
and $\OFI$ is in both the mechanical update of the mid-price of
Proposition \ref{prop.idealisedUpdate}.
Forward, the first predicts reversal and the second continuation.

The sign therefore belongs to the kernel and not to $\OFI$,
and no statement of the form ``$\OFI$ predicts continuation'' is made in these notes
without naming the regime it is made in.
No coefficient, window or goodness-of-fit figure is reproduced here.

The apparatus of Section \ref{sec.pointProcesses} generates order flow in a controlled setting,
so that a result obtained in it is a result about the generator and not about a traded market.
The following are absent by construction.

The intensities do not read the book: there is no queue-depletion feedback,
so the arrival rates are indifferent to whether a queue is about to empty.
The \emph{marks} do read it, and that is the mechanism behind the book's restoring force.
There is no queue-position economics.

There is no informed trading, and therefore no \emph{adverse selection}.
There is a single excitation timescale, the kernel being one exponential.
There is no intraday non-stationarity, so a window tuned here is not tested against one.
There is no anchor for the price level, so the level wanders and the mean price change over
a session is a property of the seed.

The first omission is the largest, and it bears on how the two statistics of this chapter
are to be compared.
In a traded market part of the flow is informed.
A resting limit order is then short an option: it is filled when someone wishes to trade
against it, which is disproportionately when trading against it is right,
and the defence available to whoever posted it is to withdraw before that happens.
A queue about to be picked off is therefore a queue whose own side pulls away from it,
and depletion begets further depletion.
That is what makes $\queueImbalance[1]$ informative in a traded book:
a thin bid is not merely closer to being cleared,
it is thin \emph{because} those posting on that side inferred something.

The generator of these notes does the reverse, by construction.
Its kernel is built so that whatever depletes a side excites the limit orders that refill it
far more than it excites further withdrawals from it,
that being the local mechanism which keeps a queue from emptying.
An adverse-selection channel is the reverse of that mechanism,
and the two cannot both be strong in one kernel.

\begin{remark}\label{remark.asymmetricReading}
 It follows that any comparison of $\queueImbalance[1]$ with $\OFI$ conducted here must be
 read asymmetrically.
 Whatever forecasting power $\queueImbalance[1]$ retains in this setting is mechanical ---
 a depleted queue is one market order from clearing, and clearing moves the mid --- and
 none of it is informational,
 so the comparison measures the imbalance with its principal advantage switched off.
 A result favourable to $\OFI$ here is therefore not evidence that $\OFI$ dominates
 $\queueImbalance[1]$ in a traded market;
 a result favourable to $\queueImbalance[1]$ would be the stronger finding.
\end{remark}

Two questions cannot be asked in this setting at all, requiring recorded data:
whether any of the regimes described here occurs in a traded market,
and whether one excitation timescale suffices to describe one.
They are the subject of Section \ref{sec.lobsterEmpirics}.
